\chapter{Heuristiques de développement}

Au fil du projet, nous nous sommes donné quelques règles. Aucune n'est originale, mais chacune vient d'une erreur que
nous avons faite au moins une fois.

\section{La fidélité d'abord}
Le moteur par défaut est toujours le plus exact dont nous disposons : reSIDfp, ares, Nuked, Ootake. Quelqu'un qui
ouvre \soft pour la première fois doit entendre le morceau tel qu'il sonnait. Les moteurs plus légers restent
disponibles pour les téléphones lents, et pour comparer.

\section{Montrer l'état exact}
Toute valeur affichée est celle de la puce à l'image en cours, sans moyenne ni lissage. La première vue PC Engine
avait un historique de la hauteur qui se recadrait en permanence pour rester lisible ; résultat, la courbe semblait
toujours en retard sur la musique. Nous l'avons remplacée par une description de la forme d'onde en mots, qui est
exacte à chaque image.

\section{Le graphique quand il suffit}
Quand une courbe dit tout, nous ne la doublons pas d'un texte. La réponse du filtre du SID affichait ses valeurs
en chiffres à côté de la courbe ; personne ne les lisait, et elles cachaient la courbe. Quand un texte est utile,
nous le mettons à côté de ce qu'il décrit, pas dans une légende à l'autre bout de l'écran.

\section{Une vue par synthèse}
Nous avons commencé par une vue générique, la même pour toutes les puces. Elle ne montrait rien d'intéressant. Chaque
synthèse a sa logique, et c'est cette logique qu'il faut montrer : la table d'ondes pour la PC Engine, les opérateurs
pour la FM, les enveloppes et l'écho pour la Super Nintendo, le filtre et les modulations pour le SID.

\section{Ne dire que ce qui est utile}
La première version de la colonne Mod du SID indiquait la source et la destination de chaque modulation, avec des
flèches. C'était exact et illisible. Comme la destination se déduit toujours de la source sur le SID, nous n'avons
gardé que la source.

\section{Mesurer avant d'affirmer}
Chaque correctif est vérifié par une mesure avant d'être livré : sorties par voix comparées à la sortie mixée, rendu
de contrôle, test. Quand nous n'avons pas pu vérifier quelque chose, nous le disons.

\section{Voir avant de livrer}
Toute modification visuelle est contrôlée sur une image produite par les tests avant d'être compilée. Cela nous a
évité de livrer plusieurs fois du texte coupé ou superposé, mais pas toujours : certains défauts, comme les marges de
l'encoche de l'écran, ne se voient que sur un vrai téléphone.

\section{Toucher le moins possible au code des autres}
Nos modifications dans les bibliothèques importées sont courtes, marquées dans le code et écrites dans le style du
fichier modifié. Le jour où ces bibliothèques évoluent, nous devons pouvoir reporter nos changements sans tout
reprendre.

\section{Ne rien imposer}
Ce qui change la disposition de l'écran (voie des samples du SID, voies de batterie de l'AdLib, mono pour les
morceaux à plusieurs SID) est désactivé par défaut. Vous l'activez si vous en avez besoin.

\section{Vérifier les licences à l'entrée}
Avant d'importer du code, nous vérifions que sa licence est compatible avec la GPL~v3. Pour certains cœurs de libvgm,
la licence n'est pas indiquée fichier par fichier ; nous les avons listés pour les vérifier avant toute publication.

\section{Pas de régression}
Avant chaque livraison importante, nous relançons toute la suite de tests. Les seuls échecs que nous acceptons sont
ceux des tests de catalogues, qui ont besoin du réseau.
